\documentclass[12pt]{article}
\usepackage[T1]{fontenc}
\usepackage[utf8]{inputenc}
\usepackage{lmodern}
\usepackage{textcomp}
\usepackage{enumitem}
\usepackage{geometry}
\geometry{margin=1in}
\begin{document}
\begin{center}
Answer Key - Economics - Graduate Level Assessment 1 \
\small (Answers are ordered exactly as in the question paper.)
\end{center}

\section*{Multiple Choice}
\begin{enumerate}[label=\arabic*.]
\item \textbf{Answer:} D
\\ \textit{Options:} A) The a, b, and c coefficients significant; B) The a and c coefficients significant; C) The a, b, and d coefficients significant; D) The b and d coefficients significant; E) The a, d, and c coefficients significant
\\ \textit{Note:} The presence of significant b and d coefficients in the VAR model indicates that changes in one variable influence another and vice versa, thus demonstrating bi-directional feedback between the variables.
\item \textbf{Answer:} A
\\ \textit{Options:} A) operates in an industry with easy entry and exit; B) will charge the highest price it can on the demand curve; C) will always make a profit in the short run; D) can increase its profit by increasing its output; E) faces a perfectly elastic demand curve
\\ \textit{Note:} The correct answer is inaccurate because a monopoly operates in an industry characterized by high barriers to entry and exit, which is fundamentally different from the conditions of perfect competition that allow for easy entry and exit.
\item \textbf{Answer:} A
\\ \textit{Options:} A) low levels of income relative to other nations.; B) a decrease in the country's production.; C) foreigners having no taste for this country's products.; D) a loose monetary policy.; E) high levels of income relative to other nations.
\\ \textit{Note:} A balance of trade surplus can occur when a country has low levels of income relative to other nations, as this can lead to reduced domestic consumption and increased exports, resulting in a higher value of goods and services sold abroad compared to what is imported.
\item \textbf{Answer:} A
\\ \textit{Options:} A) Corn would become a luxury good due to the price ceiling.; B) The government would have to subsidize corn production.; C) The price of corn would increase dramatically.; D) The price of corn would remain the same but the quality would decrease.; E) The producers of corn would lose revenue due to the decreased price.
\\ \textit{Note:} The correct answer is based on the concept that a price ceiling set above the equilibrium price does not affect the market price or quantity supplied, allowing corn to maintain its status as a normal good rather than a luxury good.
\item \textbf{Answer:} A
\\ \textit{Options:} A) monopoly output being less than the competitive output.; B) setting the price below average total cost.; C) short-run abnormal profits.; D) monopoly output being greater than the competitive output.; E) long-run abnormal profits.
\\ \textit{Note:} Monopoly deadweight loss occurs because a monopolist restricts output below the socially optimal level, leading to lost consumer and producer surplus that would otherwise be realized in a competitive market.
\end{enumerate}

\section*{True / False}
\begin{enumerate}[label=\arabic*.]
\setcounter{enumi}{5}
\item \textbf{Answer:} True
\\ \textit{Options:} A) True; B) False
\\ \textit{Note:} The statement is true because, in the long run, a monopolist maximizes profit by producing at the output level where marginal revenue equals marginal cost, which allows it to set a price that equals average total cost when it earns zero economic profit.
\item \textbf{Answer:} True
\\ \textit{Options:} A) True; B) False
\\ \textit{Note:} Decreased government spending and increased taxes can reduce overall demand in the economy, thereby helping to alleviate inflationary pressures.
\item \textbf{Answer:} False
\\ \textit{Options:} A) True; B) False
\\ \textit{Note:} Inflation causes the purchasing power of money to decrease, meaning that consumers can buy fewer goods with the same amount of money as prices rise.
\item \textbf{Answer:} False
\\ \textit{Options:} A) True; B) False
\\ \textit{Note:} A leftward shift in the aggregate supply curve typically leads to higher price levels and decreased output, which aligns with the Phillips curve's implication of a trade-off between inflation and unemployment rather than resulting in lower price levels.
\item \textbf{Answer:} True
\\ \textit{Options:} A) True; B) False
\\ \textit{Note:} Highly efficient and fast markets for stocks and bonds enable quick transactions and price adjustments, leading to greater liquidity compared to the relatively illiquid and slower market for automobiles.
\end{enumerate}

\section*{Long Form Response}
\begin{enumerate}[label=\arabic*.]
\setcounter{enumi}{10}
\item \textbf{Answer:} Communal property plays a significant role in the development and functioning of capitalist systems by providing a counterbalance to individual ownership and fostering collective welfare. This institution allows communities to pool resources and manage them collectively, which can enhance economic resilience and mitigate inequalities typically exacerbated by capitalism. The implications for economic structures include the potential for more equitable distribution of resources and the promotion of social cohesion, which can lead to increased productivity and innovation. Furthermore, communal property can challenge traditional capitalist notions of private ownership by demonstrating alternative models of cooperation and shared benefits, ultimately influencing social relations by reinforcing community bonds and shared responsibilities. Thus, while capitalism often emphasizes individualism, communal property introduces a critical dimension that can enrich both economic outcomes and social dynamics.
\\ \textit{Note:} correct because it highlights how communal property serves as a counterbalance to individual ownership in capitalist systems, promoting equity and collective welfare while addressing inequalities and enhancing economic resilience.
\item \textbf{Answer:} A major earthquake disrupts production capabilities by damaging infrastructure, displacing labor, and reducing the availability of resources, which shifts the aggregate supply curve to the left. As the aggregate supply decreases, the equilibrium quantity of output diminishes due to reduced production capacity. Simultaneously, the decrease in supply leads to increased competition for remaining goods and services, resulting in a rise in the equilibrium price level. Thus, the economy experiences a higher price level alongside a lower quantity of output, illustrating the adverse effects of supply shocks on economic equilibrium. This analysis highlights the critical interplay between supply disruptions and market dynamics in influencing both prices and output levels.
\\ \textit{Note:} correct because it accurately describes how a major earthquake disrupts production, leading to a leftward shift in the aggregate supply curve, which reduces the equilibrium quantity of output and increases the equilibrium price level due to decreased supply and heightened competition for goods and services.
\end{enumerate}

\vfill
\noindent\textit{End of Answer Key}
\end{document}